\section{Why the hub survives noise}\label{sec:robustness}

\subsection{The puzzle}
Near the lock-in edge, tracking stays at 0.91--0.99 with up to 50\% false alarms at every size. A noisy mean field
averages the reward on the source over the noise in the symmetric state (A4):
\begin{equation}
  \mathbb E[r](x)=\Big(1-\nu+\tfrac{\nu}{N}\Big)f(x)+\nu\,\tfrac{N-1}{N}\,g(x),
\end{equation}
where $f$ is the reward when the true source is seeded and $g$ when a rival is (the rival then carries the $\rho$
advantage). The averaged map $x\mapsto(x+\varepsilon\mathbb E[r])/(1+\varepsilon)$ matches the simulation within 3\%
for $T\ge0.45$ but falls far short near $\kappa$ (\cref{tab:meanfield}). Keeping the round-to-round fluctuations (a
Jensen effect) and the relay effect closes part of the gap.

\begin{table}[tbp]\centering\small
\caption{Tracking: simulation (E8) against the averaged mean field, the fluctuating walk and the relay-aware walk.}\label{tab:meanfield}
\input{generated/meanfield}
\end{table}

\subsection{Switching effects off one at a time}
Every reduction keeps the symmetric state, in which all rivals of the source hold equal attention; the simulation does
not. A $2\times2$ factorial removes that one difference: after every round, each agent's attention on its rivals is reset
to their mean, leaving the share on the source unchanged. A second switch removes relays, by rewarding direct evidence
instead of closeness (\cref{sec:learning}). Against these runs we place four reductions: the symmetric walk, the
relay-aware walk, and a column walk that keeps one agent's full unequal column with the rest of the network symmetric
(\cref{tab:a4,fig:a4}).

\figw{mechanism}{\linewidth}{Why the hub keeps its place. Left: uneven rivals pull less in total
(score $\propto w^{\kappa/T}$, concave). Middle: news from a false alarm passes through the hub. Right: tracking at
$N=10$, $\nu=0.5$, $T=0.25$, removing one effect at a time.}{mechanism}

\figpair{a4-10}{10 agents}{a4-80}{80 agents}{The A4 test at $\nu=0.5$: tracking against $T$ for each condition and
reduction.}{a4}

\begin{table}[tbp]\centering\small
\caption{The A4 test at $T=0.25$. ``Equal'': rivals equalized every round. ``Competition'': total rival pull relative to
equal rivals in the evidence column walk.}\label{tab:a4}
\resizebox{\linewidth}{!}{\input{generated/a4}}
\end{table}

\paragraph{Results.}
\begin{enumerate}
  \item With rivals equalized, each simulation matches its reduction: within 0.03 at $N\ge20$, 0.06 at $N=10$.
  \item Unequal rivals carry the whole excess near the edge: at $N=10$, $\nu=0.5$, $T=0.25$ tracking falls from $0.94$ to
    $0.75$ when rivals are equalized; at $N=80$ from $0.94$ to $0.82$. Far from the edge nothing changes.
  \item Relays amplify the effect: unequal rivals add $0.20$ with relays and $0.09$ without.
\end{enumerate}

\subsection{Mechanism: concavity}
A link's pull grows like $w^{\kappa/T}$. Above the lock-in edge $\kappa/T<1$ and the pull is concave in weight. Each false
alarm deposits extra attention on one random rival, which decays over the next few rounds; the leftover attention is
therefore uneven, and uneven weight has less total pull than the same weight spread evenly. The competition factor
$\sum_k w_k^{\kappa/T}/[(N-2)\bar w^{\kappa/T}]$ is 0.87--0.94 at $T=0.25$ and approaches one as $T$ grows --- exactly
where the excess appears and disappears. A single unequal column recovers 40--100\% of the excess (typically two
thirds); the remainder comes from relay distances, which also become unequal.

\paragraph{Relays.} When a false alarm fires at $j$, closeness is measured from $j$. Everyone listens mostly to the hub,
so the hub is the closest relay even to $j$'s news and keeps more reward than an ordinary rival would.

\paragraph{The visible trace.} With noise, treeness holds at the plateau while feedforwardness falls (0.88 to
0.57--0.78 at $T=0.25$): the short-lived side hubs form loops among the rivals (\cref{fig:ff}).

\figw{e8-ff}{0.8\linewidth}{Feedforwardness against commitment, 80 agents. Noise lowers it even where tracking is still
near one.}{ff}

\takeaway{The loops that appear under noise are how the group absorbs false alarms; equalize them away and the hub
becomes fragile.}
